%%=============================================================================
%% Conclusie
%%=============================================================================

\chapter{Conclusie}%
\label{ch:conclusie}

% TODO: Trek een duidelijke conclusie, in de vorm van een antwoord op de
% onderzoeksvra(a)g(en). Wat was jouw bijdrage aan het onderzoeksdomein en
% hoe biedt dit meerwaarde aan het vakgebied/doelgroep?
% Reflecteer kritisch over het resultaat. In Engelse teksten wordt deze sectie
% ``Discussion'' genoemd. Had je deze uitkomst verwacht? Zijn er zaken die nog
% niet duidelijk zijn?
% Heeft het onderzoek geleid tot nieuwe vragen die uitnodigen tot verder
%onderzoek?

De onderzoeksvraag was hoe een onderhoudbare ticketingapplicatie gebouwd kan worden waarin klant en administrator een duidelijke supportworkflow volgen. Het antwoord is dat onderhoudbaarheid vooral ontstaat door verantwoordelijkheden consequent te scheiden. Controllers behandelen HTTP, services voeren use cases uit, repositories schermen databanktoegang af en domeinregels bewaken de businessregels. DTO's vormen de grens tussen requestdata en applicatielogica.

Functioneel levert Forcebit Ticketing het beoogde proof of concept op. Klanten kunnen tickets aanmaken, beantwoorden, sluiten en opnieuw openen. Administrators kunnen tickets opvolgen, zoeken, filteren, beantwoorden en van status veranderen. De statusflow met \texttt{New}, \texttt{Open} en \texttt{Closed} maakt het verschil tussen een nieuwe vraag, een lopende conversatie en een afgesloten ticket duidelijk.

De frontend volgt dezelfde lijn. Schermen tonen pagina's, hooks beheren gedrag, API-modules communiceren met de backend en componenten houden herbruikbare UI apart. Daardoor konden chatweergave, afbeeldingspreview, uploadlimiet, zoekvelden, filters en paginatie toegevoegd worden zonder alle logica in de schermen te plaatsen. Documentatie legt de belangrijkste verantwoordelijkheden en uitbreidingspunten vast.

De meerwaarde voor Forcebit is dat vraag, antwoord, status en bijlagen op een plaats samenkomen in plaats van verspreid te blijven over telefoon en losse e-mails. Voor mij als student bracht het project backend, databank, authenticatie, frontend, uploads, e-mail en documentatie samen in een realistische workflow.

Kritisch bekeken blijft het resultaat een proof of concept en geen productieklare applicatie. Voor dagelijks gebruik zijn tests, logging, secret management, productiegerichte opslag en verfijnd adminbeheer nodig. Een aparte bestandsserver of object storage past beter bij langdurig gebruik. Binnen de scope toont het project dat gescheiden verantwoordelijkheden de onderhoudbaarheid haalbaar maken.
